\exercice{Surfusion du phosphore}{
On étudie un phénomène de \og retard à la solidification \fg{} : certains corps purs sont susceptibles d'exister à l'état liquide, sous une pression donnée, à une température inférieure à leur température de fusion. Ce phénomène porte le nom de surfusion. Il nécessite des conditions expérimentales particulières et peut cesser lors de l'introduction d'un cristal de solide, d'une impureté ou en cas d'agitation du récipient contenant le liquide surfondu.

Soit un récipient calorifugé contenant une masse $m_0 = 10 ~ \si{\gram}$ de phosphore liquide surfondu à la température $T_0 = 34 ~ \si{\celsius}$ sous la pression atmosphérique.

\Q On fait cesser la surfusion et on observe un nouvel état d'équilibre diphasé du phosphore. Déterminer la masse respective de chacune des phases.


On donne pour le phosphore : $T^\star_f = 317 ~ \si{\kelvin}$ ; $\Delta_{fus} h (T^\star_f ) = 20.9 \cdot 10^3 ~ \si{\joule\per\kilo\gram}$ sous la pression atmosphérique et $c_P (liq) = 0.795  ~ \si{\joule\per\gram\per\kelvin}$ .

Ces valeurs sont supposées indépendantes de la température dans l'intervalle considéré.

\Q Calculer la variation d'entropie correspondante.

\Q Quel serait l'état final du système si on faisait cesser la surfusion d'une même masse de phosphore initialement à la température $T_0' = 17.5 ~ \si{\celsius}$ ?

\textbf{Données :} $c_P (sol) = 0.840 ~ \si{\joule\per\gram\per\kelvin}$}